% Einheit 2 von 5 · Ausklammern (terme.md, Einheit 4) – Hauptnummern 16–19
\setcounter{aufgabe}{15}
\zweigkopf{Einheit 2 von 5 · Ausklammern}{Hier lernst du, einen gemeinsamen Faktor vor die Klammer zu ziehen · neu oder schon bekannt – je nach Buch · keine P10-Aufgabe · baut auf: Zahl mal Klammer (Einheit 1)}{e2}

\begin{aufgabe}{Ich kann einen gemeinsamen Faktor ausklammern. Klammere den gemeinsamen Faktor aus.}
\beispiel{4x + 20 &= 4 \cdot (x + 5)}
\begin{gleichungsraster}[1]
\gl{3x + 12} & \gl{5m + 10} \\ \noalign{\vspace{5mm}}
\gl{7y + 21} & \gl{2x - 16} \\ \noalign{\vspace{5mm}}
\gl{x^2 + 9x} & \gl{4x^2 + 12x} \\ \noalign{\vspace{5mm}}
\gl{5x^2 + 5x} & \gl{6x^2 + 3x - 9} \\ \noalign{\vspace{5mm}}
\end{gleichungsraster}
Klammere aus und mache die Probe durch Ausmultiplizieren.
\begin{gleichungsraster}[2]
\gl{10x^2 - 15x} & \gl{12m^2 + 8mn - 4m} \\ \noalign{\vspace{5mm}}
\end{gleichungsraster}
\end{aufgabe}

\begin{aufgabe}{Ich finde den Fehler beim Ausklammern. Finde den Fehler, benenne ihn und korrigiere die Rechnung. Rechne danach die Aufgabe darunter selbst.}
\begin{teile}
\teil Ben rechnet so:
\end{teile}
\rechnung{6x + 18 &= 6 \cdot (x + 18)}
\begin{gleichungsraster}[1]
\gl{8m - 24} & \\ \noalign{\vspace{5mm}}
\end{gleichungsraster}
\begin{teile}
\teil Ella rechnet so:
\end{teile}
\rechnung{7y^2 + 7y &= 7y \cdot y}
\begin{gleichungsraster}[1]
\gl{9x^2 - 9x} & \\ \noalign{\vspace{5mm}}
\end{gleichungsraster}
\end{aufgabe}

\begin{aufgabe}{Ich kann entscheiden, ob man einen Faktor ausklammern kann. Kann man einen gemeinsamen Faktor größer als 1 ausklammern? Kreuze an. Wenn ja, schreibe den größten gemeinsamen Faktor auf.}
\begin{geruest}
\gz{a}{$10x + 4$}{\janein \quad Faktor: \leerfeld}
\gz{b}{$4x + 9$}{\janein \quad Faktor: \leerfeld}
\gz{c}{$x^2 + 5$}{\janein \quad Faktor: \leerfeld}
\gz{d}{$8m^2 + 12m$}{\janein \quad Faktor: \leerfeld}
\end{geruest}
\begin{teile}
\teil Mia klammert bei $12x + 8$ so aus: $12x + 8 = 2 \cdot (6x + 4)$. Ist das falsch? Begründe.
\end{teile}
\end{aufgabe}

\begin{aufgabe}{Ich kann mit Ausklammern eine Seitenlänge bestimmen. Für ein Rechteck gilt: Flächeninhalt = Länge $\cdot$ Breite.}
\begin{teile}
\teil Ein Rechteck ist $3$~cm breit und hat den Flächeninhalt $(6k + 15)$~cm². Gib einen Term für seine Länge an.
\teil Ein anderes Rechteck ist $k$~cm breit und hat den Flächeninhalt $(k^2 + 8k)$~cm². Gib einen Term für seine Länge an.
\teil Berechne für $k = 4$ die Länge des ersten Rechtecks und prüfe mit dem Flächeninhalt.
\end{teile}
\end{aufgabe}
